% Fiche d'examen -- Integration stochastique B (CentraleSupelec MDS, cours de H. Baili)
% A4 recto-verso, modele a recopier a la main.
% RECTO = conditionnement, temps d'arret, martingales, brownien.
% VERSO = integrale d'Ito, Doob-Meyer, crochets, formule d'Ito, EDS, methode.
% Compilation : pdflatex fiche_integration_sto.tex (une passe suffit)
\documentclass[a4paper]{article}
\usepackage[utf8]{inputenc}
\usepackage[T1]{fontenc}
\usepackage{lmodern}
\usepackage[a4paper,margin=4mm,top=3.5mm,bottom=3.5mm]{geometry}
\usepackage{amsmath,amssymb,mathtools}
\usepackage{multicol}
\usepackage[dvipsnames]{xcolor}
\usepackage{microtype}
\usepackage{ragged2e}
\pagestyle{empty}
\setlength{\parindent}{0pt}
\setlength{\parskip}{0pt}
\setlength{\columnsep}{2.2mm}
\setlength{\columnseprule}{0.2pt}
\setlength{\multicolsep}{0pt}
\setlength{\abovedisplayskip}{0pt}\setlength{\belowdisplayskip}{0pt}
\setlength{\abovedisplayshortskip}{0pt}\setlength{\belowdisplayshortskip}{0pt}
\setlength{\arraycolsep}{1.5pt}
\setlength{\tabcolsep}{1.5pt}
\renewcommand{\arraystretch}{0.95}
\thickmuskip=2mu plus 2mu \medmuskip=1.5mu plus 1mu \thinmuskip=1.5mu

% ---------- macros ----------
\newcommand{\R}{\mathbb{R}}\newcommand{\N}{\mathbb{N}}\newcommand{\Q}{\mathbb{Q}}
\newcommand{\E}{\mathbb{E}}\newcommand{\PP}{\mathbb{P}}
\newcommand{\cF}{\mathcal F}\newcommand{\cG}{\mathcal G}\newcommand{\cA}{\mathcal A}\newcommand{\cB}{\mathcal B}\newcommand{\cP}{\mathcal P}
\newcommand{\Nn}{\mathcal N}
\newcommand{\ind}{\mathbf 1}
\newcommand{\dd}{\mathrm d}
\newcommand{\br}[1]{\langle #1\rangle}
\newcommand{\Var}{\operatorname{Var}}\newcommand{\Cov}{\operatorname{Cov}}
\newcommand{\nr}[1]{\|#1\|}
\newcommand{\tmax}{t_{\max}}
\newcommand{\eps}{\varepsilon}
\newcommand{\mq}{\text{m.q.}}
\newcommand{\ps}{\text{p.s.}}
\definecolor{cA}{HTML}{1F4E79}  % bleu : recto
\definecolor{cB}{HTML}{7B2C2C}  % bordeaux : verso
\definecolor{cP}{HTML}{C0392B}  % rouge : pieges
\definecolor{cK}{HTML}{1E7B34}  % vert : methode
\newcommand{\hcol}{cA}
\newcommand{\Sec}[1]{\par\vspace{1.2pt}\noindent\colorbox{\hcol}{\parbox{\dimexpr\linewidth-2\fboxsep}{\color{white}\bfseries #1}}\par\vspace{0.6pt}}
\newcommand{\B}{\par\noindent\hangindent=0.8em\hangafter=1\makebox[0.8em][l]{\textbullet}}
\newcommand{\Th}[1]{\textbf{\color{\hcol}#1}}
\newcommand{\St}{$\bigstar$}
\newcommand{\Pg}{\par\noindent\textbf{\color{cP}$\triangle$\,}}
\newcommand{\Mt}{\par\noindent\textbf{\color{cK}$\blacktriangleright$\,}}
\newcommand{\pg}[1]{{\color{cP}#1}}
\newcommand{\cq}[1]{{\color{cP}[coquille poly : #1]}}

\begin{document}
\fontsize{6.9}{7.7}\selectfont
%=====================================================================
%                              RECTO
%=====================================================================
\noindent{\fontsize{7.5}{8}\selectfont\bfseries\color{cA} RECTO -- CONDITIONNEMENT, TEMPS D'ARRÊT, MARTINGALES, BROWNIEN}\hfill{\fontsize{5}{6}\selectfont $\bigstar$=essentiel \ {\color{cP}$\triangle$=piège} \ {\color{cK}$\blacktriangleright$=méthode} \ \ $s\le t$, $\Delta W_i=W_{t_i}-W_{t_{i-1}}$}
\vspace{5pt}\begin{multicols}{3}\RaggedRight
\Sec{0. Outils d'analyse et de probabilité}
\B \textbf{Inégalités} : Markov $\PP(|X|\ge a)\le\frac{\E|X|}a$ ; Tchebychev $\PP(|X-\E X|\ge a)\le\frac{\Var X}{a^2}$ ; Cauchy--Schwarz $\E|XY|\le\nr X_2\nr Y_2$ ($\Rightarrow L^2\subset L^1$) ; Minkowski $\nr{X+Y}_2\le\nr X_2+\nr Y_2$, $\nr{\int X_t\dd t}_2\le\int\nr{X_t}_2\dd t$ ; $(a+b)^2\le2a^2+2b^2$ ; \textbf{Gronwall} $u(t)\le\alpha+\beta\int_0^tu\Rightarrow u(t)\le\alpha e^{\beta t}$ ($\alpha=0\Rightarrow u\equiv0$).
\B \textbf{Limites} : CV monotone ($0\le X_n\uparrow$) ; \textbf{CV dominée} ($|X_n|\le Y\in L^1$, $X_n\to X$ \ps\ ou en proba) ; \textbf{Fatou} $X_n\ge0$ : $\E\liminf X_n\le\liminf\E X_n$ ; \textbf{Borel--Cantelli} $\sum\PP(A_n)<\infty\Rightarrow\PP(A_n\text{ i.s.})=0$ ; Fubini ; transfert $\E g(X)=\int g\,\dd\PP_X$.
\B \textbf{Modes} : \ps\ $\Rightarrow$ proba ; \mq\ ($L^2$) $\Rightarrow$ proba $\Rightarrow$ loi ; proba $\Rightarrow$ \ps\ \pg{le long d'une sous-suite} (typewriter : \mq\ sans \ps) ; \mq\ $\Rightarrow$ CV de $\E X_n$, $\E X_n^2$ ; \textbf{$L^2(\Omega)$ complet} (Cauchy $\Rightarrow$ CV) ; \textbf{UI} + proba $\Rightarrow L^1$ ; UI si dominée ou $\sup_n\E|X_n|^p<\infty$ ($p>1$).

\Sec{1. Tribus, mesurabilité}
\B \textbf{Tribu} : $\Omega\in\cF$, stable par $A^c$ et $\bigcup$ dénombrable = information disponible. $\{\emptyset,\Omega\}$ : rien ; $\mathcal P(\Omega)$ : tout. $\sigma(X)=\{X^{-1}(B)\}$ = plus petite tribu rendant $X$ mesurable ; $X$ $\cF$-mes. $\Leftrightarrow\sigma(X)\subset\cF$.
\B \St\textbf{Cas fini} : tribu $=\sigma$(partition), atomes $A_i$ ; $n$ atomes $\Rightarrow2^n$ événements ; \textbf{$X$ $\sigma(P)$-mes. $\Leftrightarrow X$ constante sur chaque atome} ; $\sigma(X)$ : regrouper les $\omega$ de même valeur. $\cF\subset\cG\Leftrightarrow$ chaque atome de $\cG$ est dans un atome de $\cF$ ($\cG$ plus fine). Dé, $X=(0,0,5,5,10,10)$ : $\sigma(X)$ a $2^3=8$ éléments, $\{1\}\notin\sigma(X)$. Tribu pair/impair : $\ind_{\text{pair}}$ mesurable, $\ind_{\le3}$ et $\omega^2$ non.
\B \St\Th{Doob--Dynkin (TD1)} $Z$ $\sigma(\eta)$-mes. $\Leftrightarrow Z=f(\eta)$, $f$ mesurable ; d'où $\E[\xi\mid\eta]=f(\eta)$. \emph{Preuve} : $m=$ loi de $\eta$, $n(C)=\E[\ind_{\{\eta\in C\}}\xi]$ $\sigma$-add. (CV dominée), $\ll m$ ; RN sur $(Y,\mathcal Y)$ : $n(C)=\int_Cf\,\dd m=\E[\ind_{\{\eta\in C\}}f(\eta)]$ (transfert) ; or $\sigma(\eta)=\{\{\eta\in C\}:C\in\mathcal Y\}$.

\Sec{2. Espérance conditionnelle}
\B \St\textbf{Déf.} $\xi\in L^1$, $\cA\subset\cF$ : $\varphi=\E[\xi\mid\cA]$ ssi (1) $\varphi$ $\cA$-mes. ; (2) $\E[\ind_A\xi]=\E[\ind_A\varphi]\ \forall A\in\cA$ ($\Rightarrow\E[Z\xi]=\E[Z\varphi]$, $Z$ $\cA$-mes. bornée). v.a. définie \ps\ ; $\E[X\mid A]=\E[X\ind_A]/\PP(A)$ (événement) est un \emph{nombre}.
\B \textbf{Fini} : $\E[X\mid\cA]=\sum_i\ind_{A_i}\frac{\E[X\ind_{A_i}]}{\PP(A_i)}$ = moyenne locale sur chaque atome. Ex. $\Omega=\{\omega_{1..4}\}$, $\PP=(1,2,3,2)/8$, $X_3=(\frac12,\frac12,1,2)$, $\cF_1=\sigma(\{\omega_1,\omega_2\},\{\omega_3,\omega_4\})$ : $\E[X_3\mid\cF_1]=\frac12\ind_{\{\omega_1,\omega_2\}}+\frac75\ind_{\{\omega_3,\omega_4\}}$ ; tour : $\frac12\cdot\frac38+\frac75\cdot\frac58=\frac{17}{16}=\E X_3$ \checkmark.
\B \St$L^2$ : \textbf{projection orthogonale} sur $L^2(\cA)$ ($\xi-\varphi\perp L^2(\cA)$). $\E[(Y-g(X))^2]=\E[(Y-\E[Y|X])^2]+\E[(\E[Y|X]-g(X))^2]$ ; $\Var Y=\E[\Var(Y|X)]+\Var(\E[Y|X])$.
\B \Th{Unicité} $A=\{\varphi_1>\varphi_2\}\in\cA$ ; (2) : $\int_A(\varphi_1-\varphi_2)\dd\PP=0$ avec intégrande $>0$ sur $A$ $\Rightarrow\PP(A)=0$ ; symétrie.
\B \Th{Radon--Nikodym} $m$ $\sigma$-finie, $n$ signée $\sigma$-additive, $m(A)=0\Rightarrow n(A)=0$ ($n\ll m$) $\Rightarrow\exists f$ $\chi$-mes., $n(A)=\int_Af\,\dd m$.
\B \St\Th{Existence} $n(A)=\E[\ind_A\xi]$ : $\sigma$-add. ($|\sum_{i\le n}\ind_{A_i}\xi|\le|\xi|$, CV dominée), $\ll\PP$ ; RN sur \textbf{$(\Omega,\cA)$} (\pg{pas $\cF$ : la densité serait $\xi$}) $\Rightarrow\varphi$ $\cA$-mes. vérifiant (2).
\B \St\textbf{Règles} ($\cB\subset\cA$, \ps) : linéarité ; monotonie ; \textbf{tour} $\E[\E[X|\cB]]=\E X$ ; \textbf{connu reste} $X$ $\cB$-mes. $\Rightarrow\E[X|\cB]=X$ ; \textbf{connu sort} $Z$ $\cB$-mes., $XZ\in L^1$ : $\E[ZX|\cB]=Z\E[X|\cB]$ ; \textbf{petite tribu gagne} $\cB_1\subset\cB_2$ : $\E[\E[X|\cB_1]|\cB_2]=\E[\E[X|\cB_2]|\cB_1]=\E[X|\cB_1]$ ; \textbf{indép.} $\sigma(X)\perp\cB\Rightarrow\E[X|\cB]=\E X$ (\pg{réciproque fausse} : $X\sim\Nn(0,1)$, $\E[X\mid|X|]=0$). \emph{Preuve type} : vérifier (1)+(2), ex. $Z=\ind_C$ : $\E[\ind_A\ind_CX]=\E[\ind_{A\cap C}\E[X|\cB]]$, puis linéarité + CV monotone.
\B \St\textbf{Jensen cond.} $f$ convexe : $\E[f(X)|\cB]\ge f(\E[X|\cB])$ (droite d'appui en $y_0=\E[X|\cB]$ : $f(X)\ge f(y_0)+m(y_0)(X-y_0)$, $m=f'_d$). \textbf{Contraction (TD2)} $\nr{\E[X|\cG]}_2\le\nr X_2$ (Jensen $x^2$ + tour ; Pythagore) $\Rightarrow\E[\cdot|\cF_s]$ continue dans $L^2$ (sert ch.~6).
\Mt\textbf{T1 connu+indép.} $B_t=B_s+(B_t-B_s)$ : $\E[B_t|\cF_s]=B_s$, $\E[e^{aB_t}|\cF_s]=e^{aB_s+a^2(t-s)/2}$, $\E[B_t^2|\cF_s]=B_s^2+t-s$. \textbf{T2 tour} : $\E[X|B_s]=\E[\E[X|\cF_s]|B_s]$. \textbf{T3 gaussien} : $(X,Y)$ gaussien, $X-cY$ décorrélé ($c=\frac{\Cov(X,Y)}{\Var Y}$) $\Rightarrow$ indép. de $Y$ : $\E[X|Y]=\E X+\frac{\Cov(X,Y)}{\Var Y}(Y-\E Y)$ ; vect. : $\Cov(X,Y)\Cov(Y)^{-1}(Y-\E Y)$ ; $\Var(X|Y)=\Var X-\frac{\Cov(X,Y)^2}{\Var Y}$ (constante) ; loi cond. $\Nn(\E[X|Y],\Var(X|Y))$. \textbf{T4 Fubini cond.} $\E[\int_0^tX_u\dd u|\cF_s]=\int_0^t\E[X_u|\cF_s]\dd u$. \textbf{T5 geler} : $Z$ $\cB$-mes., $X\perp\cB$ : $\E[g(X,Z)|\cB]=h(Z)$, $h(z)=\E[g(X,z)]$.
\B Ex. : $X\perp Y$ : $\E[XY|Y]=Y\E X$, $\E[(X+Y)^2|Y]=\E X^2+2Y\E X+Y^2$ ; $\E[X_4|\cF_1]=\ind_{\{\omega_1,\omega_2\}}+\frac75\ind_{\{\omega_3,\omega_4\}}$ ($X_4=(1,1,1,2)$) ; $\E[\ind_A|\cB]=\PP(A|\cB)$.
\Pg Ne sortir que du $\cB$-mesurable : $\E[B_sB_t|\cF_s]=B_s\E[B_t|\cF_s]=B_s^2$ \checkmark, $\E[B_tB_s|\cF_s]=B_t\E[B_s|\cF_s]$ \pg{faux}. Indép. de $\sigma(X)$ avec \emph{toute} la tribu.

\Sec{3. Filtrations, temps d'arrêt}
\B \textbf{Filtration} $\cF_s\subset\cF_t$ ; \textbf{adapté} : $X_t$ $\cF_t$-mes. ; naturelle $\cF^X_t=\sigma(X_s,s\le t)$ ; $\cF_{t^+}=\bigcap_{s>t}\cF_s$ ; $\cF_\infty=\sigma(\bigcup\cF_t)$. \textbf{Conditions habituelles} : continue à droite ($\cF_{t^+}=\cF_t$) + complète ($\cF_0\ni$ négligeables : $Y=X$ \ps\ reste adapté).
\B Ex. $X_0=1$, $X_1=(\frac12,\frac12,2,2)$, $X_2=(1,1,1,2)$, $X_3=(\frac12,\frac12,1,2)$ : $\cF_0$ triviale, $\cF_1=\sigma(\{\omega_1,\omega_2\},\{\omega_3,\omega_4\})$, $\cF_2=\cF_3=\sigma(\{\omega_1,\omega_2\},\{\omega_3\},\{\omega_4\})$ (arbre).
\B \St\textbf{Temps d'arrêt} $\tau\in[0,\infty]$ : $\{\tau\le t\}\in\cF_t\ \forall t$ (décider sans voir le futur). \textbf{Discret} $\Leftrightarrow\{\tau=t\}\in\cF_t\ \forall t$. Filtration càd : $\Leftrightarrow\{\tau<t\}\in\cF_t$. Ex. : « vendre dès $X\ge1$ » : $\{\tau=1\}=\{\omega_3,\omega_4\}\in\cF_1$, $\{\tau=2\}=\{\omega_1,\omega_2\}\in\cF_2$ : TA ; « acheter si profit plus tard » : $\{\tilde\tau=0\}=\{\omega_3,\omega_4\}\notin\cF_0$ : \pg{pas TA}. $\tau=\ind_{\{\xi_1\ge0\}}$ : $\{\tau=0\}=\{\xi_1<0\}\notin\cF_0$.
\B \textbf{Opérations TA} : $\tau\wedge\sigma$ ($\{\tau\le t\}\cup\{\sigma\le t\}$), $\tau\vee\sigma$ ($\cap$), $\tau+\sigma$, $\tau\wedge t$, \textbf{$h(\tau)$, $h$ croissante càd, $h(t)\ge t$} ($\{h(\tau)\le t\}=\{\tau\le\tilde t\}$ ou $\{\tau<\tilde t\}$, $\tilde t=\sup\{s\le t:h(s)\le t\}\le t$) : $2\tau$, $\tau+1$, \textbf{$\tau_h=kh$ si $(k-1)h<\tau\le kh$} ($\ge\tau$, $\downarrow\tau$, dénombrable). \pg{Pas TA} : $\tau/2$, $\tau-1$ ($\{\tau\le t+1\}\in\cF_{t+1}$), dernier passage $\sup\{t\le1:W_t=0\}$, instant du max, $\arg\min$ a posteriori. \pg{Vérifier $t=0$} ($\cF_0$ souvent triviale) : ex. « 1\textsuperscript{er} instant du max de $X$ sur $\{0,..,3\}$ » : $\{\rho=0\}=\{\omega_1,\omega_2\}\notin\cF_0$.
\B \St\Th{TD3} $D$ \textbf{fermé}, $X$ adapté \textbf{continu} $\Rightarrow H_D=\inf\{t:X_t\in D\}$ TA ; $X$ \textbf{càdlàg} $\Rightarrow\tilde H_D=\inf\{t:X_t\in D\text{ ou }X_{t^-}\in D\}$ TA. \emph{Preuve} : union $\bigcup_{s\le t}$ non dénombrable ; $D=\bigcap_nD_n$, $D_n=\{d(x,D)<\frac1n\}$ ouverts ; $\{H_D\le t\}=\bigcap_n\bigcup_{r\in\Q\cap[0,t]\cup\{t\}}\{X_r\in D_n\}\in\cF_t$ ($\subset$ : continuité à droite/lim. à gauche $\Rightarrow$ un rationnel voisin tombe dans l'ouvert ; $\supset$ : $r_n\to s$, $X_s$ ou $X_{s^-}\in D$ fermé). $D$ ouvert : seulement $\{H_D<t\}\in\cF_t$ $\Rightarrow$ TA via $\cF$ càd. Ex : $H_x$, $T_1$, $T_{-x}$, explosion $\inf\{t:e^{Y_0}+W_t=0\}$, $\inf\{t:|W_t|\ge1\}$.
\B $\cF_\tau=\{A\in\cF_\infty:A\cap\{\tau\le t\}\in\cF_t\ \forall t\}$ (connu à l'instant $\tau$) ; $\tau\equiv t_0\Rightarrow\cF_{t_0}$ ; $\tau\le\sigma\Rightarrow\cF_\tau\subset\cF_\sigma$ ; $\tau$ $\cF_\tau$-mes.
\Pg $\cF_t\subset\cG_t$ : $\mathbb F$-TA $\Rightarrow\mathbb G$-TA ; \pg{réciproque fausse}.

\Sec{4. Martingales}
\B \St\textbf{Déf.} adapté, $\eta_t\in L^1$, $\E[\eta_t|\cF_s]=\eta_s$ ($s<t$) $\Leftrightarrow\E[\ind_A(\eta_t-\eta_s)]=0\ \forall A\in\cF_s$. \textbf{Sous} : $\eta_s\le\E[\eta_t|\cF_s]$ (favorable) ; \textbf{sur} : $\ge$. $\E\eta_t$ constante (croissante si sous-mart.). $\E[\xi|\cF_t]$ est une martingale (tours).
\B \St\textbf{Accroissements $\perp$ passé} : $\E[(\eta_t-\eta_s)Z]=0$, $Z\in L^2(\cF_s)$ $\Rightarrow\E\eta_t^2=\E\eta_s^2+\E(\eta_t-\eta_s)^2$, $\E[\eta_s\eta_t]=\E\eta_s^2$, $\E(M_t-M_0)^2=\E\sum_i(\Delta M_i)^2$.
\B $\Cov(\eta_s,\eta_t)=\Var(\eta_{s\wedge t})$ (comme $B$). Si $\E[\eta_t|\cF_s]=\eta_s$ pour \emph{une} filtration $\mathbb F\supset\mathbb F^\eta$, alors aussi pour $\mathbb F^\eta$ (tour + adapté). Mise prévisible $H_i$ ($\cF_{t_i}$-mes., bornée) : $\sum H_i(\eta_{t_{i+1}}-\eta_{t_i})$ martingale (« transformée ») : on ne bat pas un jeu équitable.
\B \St\Th{TD4} $f$ convexe, $f(\eta_t)\in L^1$ $\Rightarrow f(\eta)$ sous-mart. : $f(\eta_t)\ge f(\eta_s)+m(\eta_s)(\eta_t-\eta_s)$ puis $\E[\cdot|\cF_s]$. Ex. $\eta^2$, $|\eta|$, $\eta^+$.
\B \St\Th{TD6} $W_t$, $W_t^2-t$, $e^{aW_t-\frac{a^2}2t}$ martingales : $W_t^2-W_s^2=(W_t-W_s)^2+2W_s(W_t-W_s)\to\E[\cdot|\cF_s]=t-s$ ; $\E[e^{a(W_t-W_s)}]=e^{a^2(t-s)/2}$ ($\E e^Z=e^{\mu+\sigma^2/2}$). Aussi $W^3-3\int_0^tW\dd s$, $e^{t/2}\cos W_t$, $W^1W^2$ (indép.), $\exp(\int g\dd W-\frac12\int g^2)$.
\B \St\Th{Partiel 1} $\cF_t\subset\cG_t$ : $\mathbb F$-mart. $\not\Rightarrow\mathbb G$-mart. ; $\mathbb G$-mart. $\not\Rightarrow\mathbb F$-adaptée. $\cG_t=\sigma(\cF_t\cup\sigma(W_1))$ : \textbf{$W$ pas $\mathbb G$-mart.} Absurde : $s<1<t$, conditionner $\E[W_t|\cG_s]=W_s$ par $\sigma(W_1)\subset\cG_s$ : $\E[W_t|W_1]=W_1$ (accr. indép.) vs $\E[W_s|W_1]=sW_1$ (T3) $\Rightarrow W_1=sW_1$, absurde.
\B \St\Th{Arrêt 1} $\tau$ TA à valeurs $t_1<\dots<t_n$ $\Rightarrow\E\eta_\tau=\E\eta_{\tilde t}\ \forall\tilde t$. \emph{Preuve} (manches/Abel) : $\eta_\tau=\eta_{t_1}+\sum_{i<n}\ind_{\{\tau>t_i\}}(\eta_{t_{i+1}}-\eta_{t_i})$, $\ind_{\{\tau>t_i\}}=1-\ind_{\{\tau\le t_i\}}$ $\cF_{t_i}$-mes. $\Rightarrow\E[\ind\,\E[\eta_{t_{i+1}}-\eta_{t_i}|\cF_{t_i}]]=0$. \cq{lire $\eta_{t_i}-\eta_{t_{i+1}}$ sachant $\cF_{t_i}$}. Pas de domination ($|\eta_\tau|\le\sum|\eta_{t_i}|$). \Th{Arrêt 2} sous-mart., $t_n=\tmax$ : $\E\eta_\tau\le\E\eta_{\tmax}$.
\B \St\Th{Arrêt 3 (continu)} $\eta$ (sous-)mart. \textbf{càd}, $\tau$ TA \textbf{borné} $\le\tmax$, $|\eta_\tau|\le f\in L^1$ \textbf{(domination)} $\Rightarrow\E\eta_\tau=\E\eta_{\tmax}$ ($\le$). \emph{Preuve} : $\tau_h\ge\tau$ TA ; Arrêt 1 sur $\tau_h\wedge\tmax$ (valeurs finies) ; $h\to0$ : $\tau_h\wedge\tmax\downarrow\tau$, càd $\Rightarrow\eta\to\eta_\tau$ \ps, CV dominée.
\Pg $T_1=\inf\{W=1\}$ fini \ps, \textbf{non borné} : $\E W_{T_1}=1\ne0$ ; $T_1\overset{\mathcal L}=1/W_1^2$, $\PP(T_1>t)\sim\sqrt{2/(\pi t)}$, $\E T_1=+\infty$. \Mt Tronquer $\tau\wedge n$, dominer, $n\to\infty$.
\B \St\Th{Deux barrières (TD9)} $\sigma_n=T_a\wedge T_{-b}\wedge n$ borné, $|B_{\sigma_n}|\le a\vee b$, $\E B_{\sigma_n}=0$, CV dominée : $ap-b(1-p)=0$, \textbf{$\PP(T_a<T_{-b})=\frac b{a+b}$}. $I=-\inf_{s\le T_1}B_s$ : $\{I\le x\}=\{T_1\le T_{-x}\}$, $F_I=\frac x{1+x}$, \textbf{$f_I(x)=\frac1{(1+x)^2}\ind_{x\ge0}$}, $\E I=\infty$. \emph{Bonus} : $B^2-t$ $\Rightarrow\E[T_a\wedge T_{-b}]=ab$ ; $e^{\theta B-\theta^2t/2}\le e^{\theta a}$ avant $T_a$ $\Rightarrow\E e^{-\lambda T_a}=e^{-a\sqrt{2\lambda}}$.
\B \textbf{Variantes} : marche $S_n$ symétrique : $S_n$, $S_n^2-n$ mart. ; mêmes résultats ($\frac b{a+b}$, $\E\tau=ab$). Asymétrique ($p\ne q$) : $(q/p)^{S_n}$ et $S_n-(p-q)n$ mart. $\Rightarrow\PP(T_a<T_{-b})=\frac{1-(q/p)^{-b}}{(q/p)^a-(q/p)^{-b}}$. Dérive $X=W+\mu t$ : $e^{-2\mu X}$ mart. (Itô) $\Rightarrow\PP(T_a<T_{-b})=\frac{1-e^{2\mu b}}{e^{-2\mu a}-e^{2\mu b}}$. Poisson $N$ ($\lambda$) : $N_t-\lambda t$, $(N_t-\lambda t)^2-\lambda t$, $e^{\theta N_t-\lambda t(e^\theta-1)}$ mart.
\B Mart. $M\ge0$, $M_0=1$ : $\E[M_\tau\ind_{\tau<\infty}]\le1$ $\forall\tau$ (arrêt en $\tau\wedge n$ + \textbf{Fatou}, seulement $\le$) ; Ville $\PP(\sup M\ge\frac1\alpha)\le\alpha$. Arrêté $\eta^\tau$ martingale ($\eta_{n\wedge\tau}-\eta_{(n-1)\wedge\tau}=\ind_{\{\tau\ge n\}}\Delta\eta_n$, prévisible).
\B \textbf{Classe (D)} sur $[0,a]$ : $\{X_T,T\text{ TA}\le a\}$ UI : $\sup_T\E[|X_T|\ind_{|X_T|>\alpha}]\xrightarrow[\alpha\to\infty]{}0$ ; \textbf{(DL)} : (D) $\forall a$. Mart. càd $\Rightarrow$ (DL) ($M_T=\E[M_a|\cF_T]$) ; mart. $L^2$ $\Rightarrow X^2$ (DL) ($X_T^2\le\E[X_a^2|\cF_T]$). UI = permet $\lim\E=\E\lim$. \emph{Hors poly} : Doob $\PP(\sup_{s\le t}|M_s|\ge\lambda)\le\frac{\E|M_t|}\lambda$, $\E\sup_{s\le t}M_s^2\le4\E M_t^2$ ; arrêt optionnel $\E[M_\tau|\cF_\sigma]=M_\sigma$ ($\sigma\le\tau$ bornés).

\Sec{5. Mouvement brownien}
\B \St\textbf{Déf.} (1) adapté ; (2) trajectoires continues ; (3) $B_t-B_s\perp\cF_s$ ; (4) $B_t-B_s\sim\Nn(0,t-s)$ ; (5) $B_0=0$. $\Rightarrow(B_{t_1},..,B_{t_n})$ gaussien centré, \textbf{$\Cov(B_s,B_t)=s\wedge t$} ($B_t=(B_t-B_s)+B_s$). Invariances : $B_{t+t_0}-B_{t_0}$, $-B$, $c^{-1/2}B_{ct}$ (échelle), $tB_{1/t}$ (hors poly). Limite d'échelle de marches $\pm\sqrt{\delta}$ (Donsker). $B$ s'annule une infinité de fois près de $0$ ; $\limsup B=+\infty$, $\liminf=-\infty$ ; $B_t/t\to0$.
\B \textbf{Non dérivable} : $\frac{B_{t+h}-B_t}h\sim\Nn(0,\frac1{|h|})$, $\PP(|\cdot|\le c)\le2c\sqrt{|h|/2\pi}\to0$ (densité $\le$ max) : accroissement $\sim\sqrt h$.
\B \St\Th{TD5a} $\eps_{\mathbb T}=\sum(\Delta W_i)^2\to b-a$ en \mq : $\E\eps_{\mathbb T}=\sum\Delta t_i$ ; indép. : $\Var\eps_{\mathbb T}=\sum\Var(\Delta W_i)^2=2\sum\Delta t_i^2\le2(b-a)\max\Delta t_i\to0$ ($\E Z^4=3\sigma^4$, $\Var Z^2=2\sigma^4$).
\B \St\Th{TD5b} \ps, ni lipschitz ni VB sur $[a,b]$. (i) Tchebychev $\PP(|\eps_{\mathbb T_n}-(b-a)|>\delta)\le\frac{2(b-a)^22^{-n}}{\delta^2}$ (dyadiques) sommable $\Rightarrow$ Borel--Cantelli : $\Omega_0$, $\PP=1$, $\eps_{\mathbb T_n}(\omega)\to b-a$. (ii) $\omega$ fixé : \textbf{$\sum(\Delta W)^2\le\max|\Delta W|\cdot\sum|\Delta W|$} $\le c\max|\Delta W|\to0$ (unif. continuité ; lip : $\sum|\Delta W|\le K(b-a)$) : contradiction. \pg{« \ps » seulement ; \mq\ $\not\Rightarrow$ \ps : sous-suite}. $\sum|\Delta W_i|\to\infty$ ($\sim\sqrt n$). \textbf{Continu + VF $\Rightarrow$ variation quadratique nulle.}
\B \St\Th{TD7} $B_t+B_s\sim\Nn(0,s+t+2s\wedge t)$ ($=t+3s$, $s\le t$) \pg{pas $s+t$} ; $\E[B_sB_t^2]=0$ ; $\E[B_s^2B_t^2]=\E[B_s^2(B_s^2+t-s)]=$\textbf{$2s^2+st$} ($s\le t$) \cq{$3s^2+(t-s)$} ; $\E[B_t|\cF_s]=B_{s\wedge t}$ ; $\E[B_t|B_s]=B_s$ ($s\le t$), \textbf{$\frac tsB_s$} ($t\le s$ : $\Cov(B_t-cB_s,B_s)=t-cs=0$) ; $\E[\int_0^tB_u\dd u|\cF_s]=\int_0^sB_u\dd u+(t-s)B_s$ ; $\E[\int_0^tB_u\dd u|B_s]=\frac s2B_s+(t-s)B_s$.
\B Autres : $\E[W_t^2|W_s]=\frac{t^2}{s^2}W_s^2+t-\frac{t^2}s$ ($t\le s$) ; $\E[W_t^3|\cF_s]=W_s^3+3(t-s)W_s$, $\E[W_sW_t^3]=3st$ ; $\E[B_1B_2B_3]=0$ (impair) ; $B_t-2B_s\sim\Nn(0,t)$ ; $\E(W_t-W_s)^4=3(t-s)^2$ ; $\int_0^tW_s\dd s\sim\Nn(0,\frac{t^3}3)$. Pont $B_t-tB_1\perp B_1$, $\Cov=s\wedge t-st$, \pg{pas $\cF_t$-adapté}.
\B $\E[(B_t-B_s)^2|\cF_s]=t-s$ ($\br B_t=t$) ; $\E B_t^{2k}=(2k-1)!!\,t^k$ ($\E B_t^4=3t^2$, $\E B_t^6=15t^3$) ; $\E|B_t|=\sqrt{2t/\pi}$ ; $\E[e^{aB_t}]=e^{a^2t/2}$ ; $\E[B_te^{aB_t}]=at\,e^{a^2t/2}$ ; $\PP(B_t\le x)=\Phi(x/\sqrt t)$.
\B \St\Th{Markov fort} $T$ TA fini \ps\ : $(W_{T+t}-W_T)_t$ brownien pour $(\cF_{T+t})$, \textbf{indép. de $\cF_T$}. Simple (Partiel 6) : $\E[f(B_1)|\cF_t]=\E[f(B_1)|B_t]=\psi(B_t)$, $\psi(x)=\E[f(x+B_1-B_t)]$ (geler).
\B \Th{TD8} $H_x$ TA (fermé), $\tilde W=W_{H_x+\cdot}-x$ BM $\perp\cF_{H_x}$ ; $y>x$ : $H_y=H_x+\inf\{t:\tilde W_t=y-x\}$ $\Rightarrow H_y-H_x\perp\cF_{H_x}$, $\sim H_{y-x}$ : accr. indép. stationnaires. Pas BM (croissant), \textbf{sauts} aux max locaux (subordinateur).
\B \textbf{Réflexion} : $\PP(\sup_{s\le t}W_s\ge a)=2\PP(W_t\ge a)=\PP(|W_t|\ge a)$ ; $T_a\overset{\mathcal L}=a^2/W_1^2$, densité $\frac{a}{\sqrt{2\pi t^3}}e^{-a^2/2t}$.
\B \St\textbf{Gaussiennes} $Z\sim\Nn(\mu,\sigma^2)$ : $\E e^{\lambda Z}=e^{\lambda\mu+\lambda^2\sigma^2/2}$ ; centrée : $\E Z^2=\sigma^2$, $\E Z^4=3\sigma^4$, impairs $=0$, $\E|Z|=\sigma\sqrt{2/\pi}$, densité $\le\frac1{\sqrt{2\pi\sigma^2}}$. Vecteur gaussien : \textbf{décorrélé $\Leftrightarrow$ indép.} ; lim. \mq/loi de gaussiennes = gaussienne (fct. caract.).

\Sec{Fil rouge (questions de cours)}
\B Information = tribu, conditionner = projeter. Non-anticipation (adapté, TA, point gauche, prévisible) $\Rightarrow\E[Y_i\Delta W_i]=0$ : centrage, isométrie, martingale. Martingale = jeu équitable, imbattable par un TA borné/dominé. $[W]_t=t$, $(\dd W)^2=\dd t$ $\Rightarrow$ le $\frac12f''$ d'Itô et le $-\frac12(b-a)$. Intégrale = isométrie sur une classe simple + complétude de $L^2$ (deux fois : $\dd W$, $\dd M$ avec $\dd t\to\dd\br M$). Mart. continue VF = constante $\Rightarrow$ $\dd t$ et $\dd W$ séparés, décomposition unique, table $\dd t\,\dd W=0$.
\end{multicols}

%=====================================================================
%                              VERSO
%=====================================================================
\newpage
\renewcommand{\hcol}{cB}
\noindent{\fontsize{7.5}{8}\selectfont\bfseries\color{cB} VERSO -- INTÉGRALE D'ITÔ, DOOB--MEYER, CROCHETS, FORMULE D'ITÔ, EDS}\hfill{\fontsize{5}{6}\selectfont $\dd X=a\,\dd t+b\,\dd W$ ; $(Z\cdot X)_t=\int_0^tZ\,\dd X$ ; \mq\ = moyenne quadratique}
\vspace{5pt}\begin{multicols}{3}\RaggedRight
\Sec{6. Intégrale d'Itô $\int f\,\dd W$}
\B $W$ ni dérivable ni VB $\Rightarrow$ pas de Stieltjes $\omega$ par $\omega$ : on définit $\int f\dd W$ dans \textbf{$L^2(\Omega)$} par isométrie.
\B \St\textbf{Étape 1} : $f=\sum_iY_i\ind_{]t_{i-1},t_i]}$, \textbf{$Y_i$ $\cF_{t_{i-1}}$-mes.} (non anticipatif), $\E Y_i^2<\infty$ : $\int f\dd W=\sum_iY_i\Delta W_i$. \textbf{(3)} $\E\int f\dd W=0$ ($Y_i\perp\Delta W_i$). \textbf{(4) Isométrie} $\E[(\int_a^bf\dd W)^2]=\int_a^b\E[f^2]\dd t$ : diag. $\E[Y_i^2\Delta W_i^2]=\E[Y_i^2]\Delta t_i$ ; croisés $i<j$ : $Y_i\Delta W_iY_j$ $\cF_{t_{j-1}}$-mes. $\perp\Delta W_j$ $\Rightarrow0$.
\B \St\textbf{Étape 2} : $f$ adaptée, $L^2$, \textbf{unif. continue en \mq} ($|s-t|<\delta\Rightarrow\E(f_s-f_t)^2<\eps$). $f_{\mathbb T}=\sum f(t_{i-1})\ind_{]t_{i-1},t_i]}$ (\textbf{point gauche}). $\E[(\int f_{\mathbb T}\dd W-\int f_{\mathbb T'}\dd W)^2]=\int\E[(f_{\mathbb T}-f_{\mathbb T'})^2]\dd t<\eps$ (extrémités gauches à $<\delta$) $\Rightarrow$ \textbf{Cauchy dans $L^2(\Omega)$ complet} $\Rightarrow$ limite $=:\int f\dd W$ (indép. de la suite : mélanger). (3),(4) passent à la limite ; linéaire ; Chasles. Général : $f$ progressive, $\E\int f^2<\infty$ (classe monotone) ; $\int f^2<\infty$ \ps\ seul. $\Rightarrow$ \pg{mart. locale} ($T_n=\inf\{\int_0^tf^2\ge n\}$).
\B \St\textbf{$X_t=\int_0^tf\dd W$ martingale} : $\E X_t^2<\infty\Rightarrow L^1$ ; adapté (lim. $L^2$, complétude) ; escalier ($s=t_m$, $i>m$) : $\E[Y_i\Delta W_i|\cF_s]=\E[Y_i\E[\Delta W_i|\cF_{t_{i-1}}]|\cF_s]=0$ ; limite : $\nr{\E[\int_s^tf_{\mathbb T}\dd W|\cF_s]-\E[\int_s^tf\dd W|\cF_s]}_2\le\nr{\int_s^t(f_{\mathbb T}-f)\dd W}_2\to0$ (TD2). $\br X_t=\int_0^tf^2\dd s$ ; $\E[(\int_s^tf\dd W)^2|\cF_s]=\E[\int_s^tf^2|\cF_s]$ ; $\E[\int f\dd W\int g\dd W]=\E\int fg\,\dd t$.
\B \St\Th{TD5c} \textbf{$\int_a^bW\dd W=\frac12(W_b^2-W_a^2)-\frac12(b-a)$}. \emph{Preuve} : $(W_b-W_a)^2=\sum\Delta W_i^2+2\sum_i(W_{t_{i-1}}-W_a)\Delta W_i$ $\Rightarrow2\sum W_{t_{i-1}}\Delta W_i=(W_b-W_a)^2-\eps_{\mathbb T}+2W_a(W_b-W_a)\xrightarrow{\mq}W_b^2-W_a^2-(b-a)$. Point droit : $+\frac12(b-a)$ ; milieu (Stratonovich) : $\frac12(W_b^2-W_a^2)$, \pg{pas martingale}. $\E[(\int_0^1W\dd W)^2]=\int_0^1s\,\dd s=\frac14(3-2+1)=\frac12$ \checkmark.
\B \St\textbf{Wiener} ($g$ déterministe $L^2$) : $\int_a^bg\dd W\sim\Nn(0,\int_a^bg^2)$, $\perp\cF_a$ ; $Y_t=\int_0^tg\dd W$ processus gaussien, $\Cov(Y_s,Y_t)=\int_0^{s\wedge t}g^2$ ; $\E e^{\int g\dd W}=e^{\frac12\int g^2}$. (Sommes $\sum g(t_{i-1})\Delta W_i$ gaussiennes, lim. \mq) Cond. : $\E[Y_t|\cF_s]=Y_s$, $\Var[Y_t|\cF_s]=\int_s^tg^2$ ; $I|\cF_t\sim\Nn(m,v)\Rightarrow\E[e^{-I}|\cF_t]=e^{-m+v/2}$.
\B $\E[\int_0^tf\dd W|\cF_s]=\int_0^sf\dd W$ ; $\Cov(\int_0^sf\dd W,\int_0^tf\dd W)=\E\int_0^{s\wedge t}f^2$ ; $W_t$ est unif. continu en \mq\ ($\E(W_t-W_s)^2=|t-s|$). \pg{Sans non-anticipation}, $Y_i=\mathrm{sgn}(\Delta W_i)$ gagnerait à coup sûr.
\B Ex. : $\E\int_0^tW^2\dd W=0$, $\E[(\int_0^tW^2\dd W)^2]=\int_0^t3s^2=t^3$ ; $\int_0^te^s\dd W\sim\Nn(0,\frac{e^{2t}-1}2)$ ; $\int_0^ts\,\dd W=tW_t-\int_0^tW\dd s$. \pg{$f=W_{t+1}$ anticipatif} : $\E\sum W_{t_{i-1}+1}\Delta W_i=1\ne0$.

\Sec{7. Doob--Meyer, prévisible, $\br M$, $\int f\,\dd M$}
\B \textbf{Croissant} : $A_0=0$, croissant, càd, $A_t\in L^1$. \St\Th{Doob--Meyer} $X$ sous-mart. càd : $X=M+A$ ($M$ mart. càd, $A$ croissant) \textbf{$\Leftrightarrow X$ de classe (DL)} ; alors \textbf{$\exists!$ décomposition avec $A$ prévisible} (compensateur). Discret : $A_n=\sum_{k\le n}\E[X_k-X_{k-1}|\cF_{k-1}]$ ($\ge0$, $\cF_{n-1}$-mes.), $M=X-A$. Ex. $W^2=(W^2-t)+t$ ; Poisson : $N=0+N$ et $N=(N-\lambda t)+\lambda t$ : \pg{unicité seulement avec $A$ prévisible} ($\lambda t$).
\B \textbf{Tribu prévisible} $\cP$ sur $\R_+\times\Omega$ : engendrée par les processus adaptés \textbf{continus à gauche} = par les rectangles $\{0\}\times A$ ($A\in\cF_0$), $]s,t]\times A$ ($A\in\cF_s$). « Connu juste avant ». Adapté continu $\Rightarrow$ prévisible.
\B \St\textbf{$\br X$} ($X$ mart. càd $L^2$) = compensateur de $X^2$ (sous-mart. (DL)). \textbf{$\E[(X_t-X_s)^2|\cF_s]=\E[X_t^2|\cF_s]-X_s^2=\E[\br X_t-\br X_s|\cF_s]$} (développer, sortir $X_s$, $X^2=M+A$). $\br W_t=t$ ; $\br{\int f\dd W}=\int f^2\dd s$ (déterministe si $f$ l'est).
\B \St\textbf{$\int\dd M$} : $\dd\br M$ mesure de Stieltjes ($\omega$ par $\omega$) ; $L^2(\cP,\mu_{\br M}\times\PP)$ : $f$ prévisible, $\nr f^2=\E\int f^2\dd\br M<\infty$. Simple $f=\sum\alpha_i\ind_{A_i}\ind_{]t_{i-1},t_i]}$, $A_i\in\cF_{t_{i-1}}$ : $I(f)=\sum\alpha_i\ind_{A_i}\Delta M_i$. \textbf{Isométrie} $\E I(f)^2=\nr f^2$ : croisés $\E[\cdots\E(\Delta M_j|\cF_{t_{j-1}})]=0$ (propriété de martingale, pas la disjonction des $A_i$ du poly) ; diag. $\alpha_i^2\E[\ind_{A_i}\E(\Delta M_i^2|\cF_{t_{i-1}})]=\alpha_i^2\E[\ind_{A_i}\Delta\br M_i]$. \textbf{Prolongement} : simples denses (classe monotone), $\nr{I(f_n)-I(f_m)}_{L^2(\Omega)}=\nr{f_n-f_m}$ Cauchy $\Rightarrow I(f)=\lim I(f_n)$ ; $\E I=0$, $\E I^2=\nr f^2$, $\br{\int f\dd M}=\int f^2\dd\br M$. Dictionnaire : $W\leftrightarrow M$, $\dd t\leftrightarrow\dd\br M$, continuité unif. \mq\ $\leftrightarrow$ classe monotone.
\Pg $\br M=[M]$ si $M$ continue ; Poisson compensé $M=N-\lambda t$ : $\br M=\lambda t$, $[M]=N$.

\Sec{8. Martingales locales, semimart., crochets}
\B Arrêté $Y^T_t=Y_{t\wedge T}$ ($0$ sur $\{T=0\}$). \St\textbf{Mart. locale} : $\exists T_n\uparrow\infty$ TA, $X^{T_n}$ martingale. Mart. $\Rightarrow$ locale ($T_n=n$) ; locale bornée/dominée $\Rightarrow$ vraie ; locale $\ge0\Rightarrow$ sur-mart. (Fatou). Utilité : pas d'intégrabilité ; processus définis sur $[0,S[$ (explosion). \pg{locale $\not\Rightarrow$ intégrable}.
\B \textbf{Semimart.} $X=X_0+M+A$, $M$ mart. locale càdlàg, $A$ adapté \textbf{VF} (= différence de 2 croissants, $\sum|\Delta A|\le c$). $X$ continue : $M,A$ continus, \textbf{décomposition unique}.
\B \St\Th{TD10} $M$ mart. $L^2$ continue VF $\Rightarrow$ constante : $\E(M_t-M_0)^2=\E\sum\Delta M_i^2$ (orthog.) $\le\E[c\max|\Delta M_i|]\to0$ (unif. cont., dominé par $c^2$) ; localiser ($T_n$ + arrêt quand variation $>k$). Unicité : $M_1-M_2=A_2-A_1$ mart. loc. continue VF $=0$.
\B \St\textbf{Covariation} $[X,Y]_t=X_tY_t-X_0Y_0-\int_0^tX_{s^-}\dd Y_s-\int_0^tY_{s^-}\dd X_s$ (= IPP). \Th{Th.} $\sum\Delta X_i\Delta Y_i\xrightarrow{\PP}[X,Y]_t$ (pas $\to0$). $[X]$ croissant ; $4[X,Y]=[X+Y]-[X-Y]$ $\Rightarrow$ VF ; $[B]_t=t$ ; $[W^1,W^2]=0$ indép. (somme centrée de var. $\sum\Delta t_i^2$) ; \textbf{$[X,A]=0$, $A$ continu VF}.
\B \St\Th{TD11} $[B+A]=[B]=t$ : $\sum(\Delta B+\Delta A)^2$, termes $\le\max|\Delta B|\sum|\Delta A|$, $\max|\Delta A|\sum|\Delta A|\to0$. D'où \textbf{$[X]=[M]$} (ignore la dérive).
\B \textbf{L1} $[Z\cdot X,Y]=Z\cdot[X,Y]$ ; \textbf{L2} $[X\cdot M,Y\cdot N]=XY\cdot[M,N]$ ($X,Y$ prévisibles loc. bornés). $M$ mart. loc. continue : $M^2-[M]$ mart. loc. ; $[M]\equiv0\Rightarrow M$ constante ; $\E M_t^2=\E[M]_t$ ($M_0=0$).
\B \St\textbf{Table} : $\dd t\,\dd t=\dd t\,\dd W=0$, $\dd W\dd W=\dd t$, $\dd W\dd W'=0$ (indép.), $=\rho\,\dd t$ si corrélés. \textbf{$\dd[X,Y]=b\,e\,\dd t$} pour $\dd Y=c\,\dd t+e\,\dd W$, \textbf{$\dd[X]=b^2\dd t$}.

\Sec{9. Formule d'Itô}
\B \St$f\in C^2(\R^d)$, $X$ semimart. continues : $f(X_t)=f(X_0)+\sum_i\int_0^tD_if(X)\dd X^i+\frac12\sum_{i,j}\int_0^tD_{ij}f(X)\dd[X^i,X^j]$. 1D : \textbf{$\dd f(X)=f'(X)\dd X+\frac12f''(X)\dd[X]$}. $f(t,x)\in C^{1,2}$ :
\[\dd f(t,X_t)=\partial_tf\,\dd t+\partial_xf\,\dd X+\tfrac12\partial_{xx}f\,\dd[X]\]
\[\qquad=\big(\partial_tf+a\,\partial_xf+\tfrac12b^2\partial_{xx}f\big)\dd t+b\,\partial_xf\,\dd W.\]
Générateur $L=a\partial_x+\frac12b^2\partial_{xx}$ ; $\frac{\dd}{\dd t}\E f(X_t)=\E[Lf(X_t)]$. \St\textbf{IPP} $\dd(XY)=X\dd Y+Y\dd X+\dd[X,Y]$ (Itô pour $xy$).
\B \emph{Preuve} : vraie pour $x$ ; vraie pour $f,g\Rightarrow fg$ (IPP sur $f(X),g(X)$ ; le terme croisé $2f'g'$ vient de $[f(X),g(X)]=\int f'g'\dd[X]$ par L2 et $[X]=[M]$) $\Rightarrow$ polynômes ; $C^2$ : Weierstrass ($f_n,f_n',f_n''\to$ unif. sur compacts), localisation, isométrie. \emph{Heuristique} : Taylor ordre 2, $(\dd X)^2=\dd[X]\sim\dd t$, $(\dd X)^3\sim\dd t^{3/2}\to0$.
\B \St\textbf{Exemples} : $\dd(W^2)=2W\dd W+\dd t$ ; $\dd W^3=3W^2\dd W+3W\dd t$ ; $\dd e^{aW-\frac{a^2}2t}=a\,e^{aW-\frac{a^2}2t}\dd W$ ; $\dd(e^{t/2}\cos W)=-e^{t/2}\sin W\dd W$ ; $\dd(tW)=W\dd t+t\dd W$ ; GBM $X=X_0e^{(\mu-\frac{\sigma^2}2)t+\sigma W}$ : $\dd X=\mu X\dd t+\sigma X\dd W$, $\dd\ln X=(\mu-\frac{\sigma^2}2)\dd t+\sigma\dd W$ ; $X^\alpha$ mart. loc. ssi $\alpha=1-\frac{2\mu}{\sigma^2}$ (ou $0$) ; $F=e^{-cW+\frac{c^2}2t}$ : $\dd F=c^2F\dd t-cF\dd W$.
\B \textbf{Moments par Itô} : $\dd W^n=nW^{n-1}\dd W+\frac{n(n-1)}2W^{n-2}\dd t\Rightarrow\E W_t^n=\frac{n(n-1)}2\int_0^t\E W_s^{n-2}\dd s$ ($\E W^4=6\int s=3t^2$). \textbf{Exponentielle} $\mathcal E_t=\exp(\int_0^t\theta\dd W-\frac12\int_0^t\theta^2\dd s)$ : $\dd\mathcal E=\theta\,\mathcal E\,\dd W$ (mart. loc. $>0$, sur-mart.). $\dd(W^1W^2)=W^1\dd W^2+W^2\dd W^1$ (indép. : martingale). $\dd X=X\dd W^1$, $\dd Y=Y(\rho\dd W^1+\sqrt{1-\rho^2}\dd W^2)$ : $\dd[X,Y]=\rho XY\dd t$.
\Mt\textbf{Recette} : (1) $f$ et le vecteur $(t,X)$, $f\in C^{1,2}$ ? (2) $\partial_t,\partial_x,\partial_{xx}$ ; (3) crochets par la table ; (4) assembler, regrouper $\dd t$ / $\dd W$ ; (5) exploiter : intégrer ; $\E$ (terme $\dd W$ centré si $\E\int(\cdot)^2<\infty$) $\Rightarrow$ EDO ; $\E[\cdot|\cF_s]$ ; \textbf{pas de $\dd t$ $\Rightarrow$ mart. locale} (+ intégrabilité $\Rightarrow$ mart.).
\Pg Oublier $\frac12f''\dd[X]$, le $\frac12$, $\partial_tf\dd t$, $\dd[X,Y]$ (nul si un facteur est VF) ; $\dd[X]=$ (coef. de $\dd W$)$^2$ ; $f\notin C^2$ ($|x|$).

\Sec{10. EDS}
\B $\dd X=f(t,X)\dd t+g(t,X)\dd W$, $X_{t_0}=X_0$ \textbf{$\Leftrightarrow X_t=X_0+\int_{t_0}^tf(s,X_s)\dd s+\int_{t_0}^tg(s,X_s)\dd W_s$} (Riemann + Itô) ; solution = semimart. continue, $[X]_t=\int g^2\dd s$.
\B \St\Th{Existence} $f,g$ continues, \textbf{$c$-lipschitz en $x$}, $X_0\in L^2$ indép. du futur de $W$ $\Rightarrow\exists$ solution adaptée, $L^2$, continue en \mq, trajectoires continues \ps\ \textbf{Picard} : $X^{(0)}=X_0$, $X^{(n)}=X_0+\int f(s,X^{(n-1)})\dd s+\int g(s,X^{(n-1)})\dd W$. Outils : Minkowski $\nr{\int X\dd t}\le\int\nr X\dd t$, isométrie, croissance linéaire $|f|\le K+c|x|$ ($f(t,0)$ bornée sur compact). Récurrence : \textbf{$\nr{X^{(n)}_t-X^{(n-1)}_t}_2\le c_1c_2^{n-1}\frac{(t-t_0)^{n/2}}{\sqrt{n!}}$}, $c_1=(K+c|x_0|)(1+\sqrt T)$, $c_2=c(1+\sqrt T)$ [$\int_0^u\frac{s^{n-1}}{(n-1)!}=\frac{u^n}{n!}$ ; $(t-t_0)^{\frac{n+1}2}\le\sqrt T(t-t_0)^{\frac n2}$ ; $2\sqrt n\le n+1$]. Série CV ($n!$ l'emporte), Cauchy unif. en $t$ ; sous-suite \ps\ unif. $\Rightarrow$ continuité ; limite dans l'équation (Lipschitz + isométrie). \textbf{Unicité} : $\E(X_t-Y_t)^2\le2c^2(T+1)\int\E(X-Y)^2$ + Gronwall. Lipschitz local $\Rightarrow$ jusqu'à explosion.
\B \St\textbf{A. Bruit additif (TD13)} $\dd X=bX\dd t+\dd B$ : $Y=e^{-bt}X$, $\dd Y=e^{-bt}\dd B$ ($\partial_{xx}=0$), \textbf{$X_t=e^{bt}(x_0+\int_0^te^{-bs}\dd B_s)$}. $\E Y_t=x_0$, $\E Y_t^2=x_0^2+\frac{1-e^{-2bt}}{2b}$ ; $Y$ gaussien $\Rightarrow\int_0^tY_s\dd s$ gaussien (Riemann, comb. lin., lim. \ps $\Rightarrow$ loi) ; $Y_t=Y_s+\int_s^te^{-bu}\dd B_u$ : $\E[Y_t|\cF_s]=Y_s$, $\Var=\int_s^te^{-2bu}\dd u$ ; \textbf{$\E[X_t|\cF_s]=e^{b(t-s)}X_s$, $\Var[X_t|\cF_s]=\frac{e^{2b(t-s)}-1}{2b}$} ; mart. ssi $b=0$. Général $\dd X=b(t)X\dd t+\sigma(t)\dd B$ : facteur $e^{-\int b}$. \textbf{OU} $\dd X=-\alpha X\dd t+\sigma\dd W$ : $X_t\sim\Nn(x_0e^{-\alpha t},\frac{\sigma^2}{2\alpha}(1-e^{-2\alpha t}))\to\Nn(0,\frac{\sigma^2}{2\alpha})$ ; stationnaire : $\Cov=\frac{\sigma^2}{2\alpha}e^{-\alpha|t-s|}$.
\B \St\textbf{B. Bruit multiplicatif (TD12)} $\dd X=m(X)\dd t+cX\dd W$ : $F=e^{-cW+\frac{c^2}2t}$, $\dd F=c^2F\dd t-cF\dd W$, \textbf{$\dd[F,X]=-c^2FX\dd t$} $\Rightarrow\dd(FX)=F\,m(X)\dd t$ : $Y=FX$, $\dot Y=F\,m(Y/F)$ EDO $\omega$ par $\omega$, $X=Y/F$. $m=\frac1x$ : $(Y^2)'=2F^2$, \textbf{$X_t=\frac1{F_t}\sqrt{X_0^2+2\int_0^tF_s^2\dd s}$} (signe $+$ : $Y^2\ge X_0^2>0$, continuité). $m=\mu x$ : GBM ; $m=0$ : $X_0e^{cW-c^2t/2}$.
\B \St\textbf{C. Changement de variable (Partiel 4)} $X=u(Y)$, $\dd Y=a(Y)\dd t+b(Y)\dd W$ $\to\dd X=(a_1X+a_2)\dd t+(b_1X+b_2)\dd W$ : \textbf{$u'a+\frac12u''b^2=a_1u+a_2$, $u'b=b_1u+b_2$}. $a=-\frac12e^{-2y}$, $b=e^{-y}$, $b_1=0$, $b_2=1$ : $u=e^y$, $a_1=a_2=0$, $X=e^{Y_0}+W$, \textbf{$Y_t=\ln(e^{Y_0}+W_t)$} sur $[0,T[$, $T=\inf\{t:e^{Y_0}+W_t=0\}$ TA (explosion). $Y$ sur-mart. ($\ln$ concave).
\B \St\textbf{D. Moments (Partiel 3)} $\dd X=\mu(t)X\dd t+\sigma(t)X\dd W$ : $\dd(X^2)=(2\mu+\sigma^2)X^2\dd t+2\sigma X^2\dd W$, $u=\E X^2$ : $u'=(2\mu+\sigma^2)u$, \textbf{$u=\E X_0^2\,e^{\int_0^t(2\mu+\sigma^2)}$} \cq{$\int2\mu\sigma^2$} ; \textbf{$\to0$ en \mq\ ssi $2\mu+\sigma^2<0$} (\pg{$\mu<0$ ne suffit pas}). GBM : $\E X_t=X_0e^{\mu t}$, $\E X_t^2=X_0^2e^{(2\mu+\sigma^2)t}$, $\Var=X_0^2e^{2\mu t}(e^{\sigma^2t}-1)$.
\B \St\textbf{EDS linéaire générale} $\dd X=(a_1X+a_2)\dd t+(b_1X+b_2)\dd W$ : $\Phi_t=e^{(a_1-\frac{b_1^2}2)t+b_1W_t}$ ($\dd\Phi=a_1\Phi\dd t+b_1\Phi\dd W$), $\dd(\Phi^{-1}X)=\Phi^{-1}(a_2-b_1b_2)\dd t+\Phi^{-1}b_2\dd W$ (crochet !) : \textbf{$X_t=\Phi_t\big(X_0+\int_0^t\Phi_s^{-1}(a_2-b_1b_2)\dd s+\int_0^t\Phi_s^{-1}b_2\dd W_s\big)$}. Moments : $m=\E X$ : $m'=a_1m+a_2$ ; $v=\E X^2$ : $v'=(2a_1+b_1^2)v+2(a_2+b_1b_2)m+b_2^2$. OU par D : $v'=-2\alpha v+\sigma^2$, $v_\infty=\frac{\sigma^2}{2\alpha}$.
\B \St\textbf{E. Taux court (Partiel 5)} $\dd r=\theta\dd t+\sigma\dd W$ : $\int_t^Tr=r(t)(T-t)+\frac\theta2(T-t)^2+\sigma I'$, $I'=\int_t^T(W_s-W_t)\dd s=\int_t^T(T-s)\dd W_s\sim\Nn(0,\frac{(T-t)^3}3)$, $\perp\cF_t$ ($\iint(u\wedge v-t)=\frac{L^3}3$) ; \textbf{$p(t,T)=\exp\{-r(t)(T-t)-\frac\theta2(T-t)^2+\frac{\sigma^2}6(T-t)^3\}$}. \pg{$\int W\dd s$ = Riemann, pas d'isométrie.}

\Sec{11. Méthode, réflexes, pièges}
\Mt\textbf{(a) TA / martingale} : définition + bonne filtration ; temps d'atteinte d'un fermé ; $\E[\eta_t-\eta_s|\cF_s]=0$ par connu+indép. \textbf{(b) calcul} : sortir/moyenner, tour, décorréler, Fubini, Wiener $\Nn(0,\int g^2)$. \textbf{(c) EDS} : additif $e^{-\int b}$ ; $cX\dd W$ $\to F$ ; $u(Y)$ identifier ; moments : Itô + $\E$. \textbf{(d) structure} : $\E(M_t-M_0)^2=\E\sum\Delta M^2$ et $\sum\Delta X^2\le\max|\Delta X|\sum|\Delta X|$, puis CV dominée.
\Mt\textbf{Preuves à savoir refaire} : unicité/existence de $\E[\cdot|\cA]$ ; arrêts 1 et 3 ; TD3 ; isométrie (étape 1) + Cauchy ; $\int W\dd W$ ; TD5 (\mq\ puis \ps) ; $\int f\dd W$ martingale ; lemme $\br X$ ; TD10, TD11 ; Itô pour un produit ; estimation de Picard.
\Mt\textbf{Vérifier} : hypothèses (TA borné + domination ; $\E\int f^2<\infty$ ; $C^2$ ; Lipschitz) ; mesurabilité du résultat ; $t=s$ redonne la valeur connue ; $\sigma=0$ redonne l'EDO ; variance $\ge0$.
\Pg Sortir $B_t$ ($t>s$) de $\E[\cdot|\cF_s]$ ; confondre $\E[B_t|\cF_s]$ et $\E[B_t|B_s]$ ; dernier passage pris pour un TA ; arrêt sans troncature ; $\int W\dd W=\frac12W^2$ ; point droit ; espérance d'une mart. locale ; martingale sans filtration ; $\br M$ vs $[M]$ avec sauts ; décomposition unique avec sauts.
\end{multicols}
\end{document}
